\documentclass[10pt]{article}
\usepackage[margin=20mm]{geometry}
\usepackage{amsmath,amssymb,mathtools,fontspec,unicode-math,xurl,hyperref}
\IfFontExistsTF{DejaVu Serif}{\setmainfont{DejaVu Serif}}{\setmainfont{Latin Modern Roman}}
\setmathfont{latinmodern-math.otf}
\setlength{\parindent}{0pt}\setlength{\parskip}{6pt}\sloppy
\title{Initial mathematical review: M25-0160
A Schiffer-type problem for annuli with applications to stationary planar Euler flows}
\author{}
\date{26 October 2026}
\begin{document}
\maketitle
Authors: Alberto Enciso, Antonio J. Fernández, David Ruiz, and Pieralberto Sicbaldi. Publication identity: Duke Mathematical Journal 174 (2025), no. 6, 1151–1208; DOI 10.1215/00127094-2024-0054 (https://doi.org/10.1215/00127094-2024-0054).\par
Reviewed version: arXiv:2309.07977v2, dated 12 August 2024, 50 pages, with its matching deposited TeX and two figures. All page numbers below refer to this author manuscript; printed and physical PDF page numbers coincide. The publisher version of record was not acquired, so the findings are not attributed to its text.\par
Initial assessment: Major, at the level of two stated proof derivations. The passage from a local Euler solution to the global weak equation drops a nonzero pressure term on the hole and the compensating boundary flux. An explicit admissible test field makes the displayed intermediate integral nonzero. In the infinite-spectrum rigidity argument, the vector used after (6.10) omits the absolute inverse-square-root weights; the asserted linear rewriting is false. Both steps directly prove advertised contributions, Theorem 1.3 and Theorem 6.1, respectively. Both admit explicit repairs verified below. No counterexample to either repaired theorem, or to the principal large-angular-mode bifurcation construction, was established. The scientific decision remains `pending`; `human\_reviews=[]`.\par
\section*{Scope and logical structure}
The main construction seeks a nonconstant Neumann eigenfunction satisfying \[\Delta u+\mu u=0\quad\hbox{in }\Omega,\qquad
 \nabla u=0\quad\hbox{on }\partial\Omega.\] Its boundary value may differ between connected components. This distinction from the classical Schiffer problem is essential. The domains are small analytic perturbations of an annulus, and the nonradial conclusion concerns the nontrivial branch, with the radial annulus retained at the bifurcation parameter zero.\par
The proof first identifies a crossing of the second nonconstant radial Neumann eigenvalue \(\mu_{0,2}\) with the first angular-mode-\(l\) Dirichlet eigenvalue \(\lambda_{l,0}\). It proves nonresonance in a dihedral symmetry class, establishes transversality for sufficiently large \(l\), and applies the Crandall–Rabinowitz theorem in anisotropic Hölder spaces. A separate numerical appendix treats \(l=4\). The resulting eigenfunction yields a Pompeiu-type integral identity and a compactly supported stationary Euler flow. Section 6 gives two additional rigidity theorems.\par
The primary reviewer read the introduction, the crossing and nonresonance argument, the complete applications section, and the decisive curvature passage. Three actual native scientific assistants checked the full spectral section; the functional setting, bifurcation proof and pullback appendix; and the rigidity section with the numerical appendix. Thus reading coverage is collective and the proof checks are bounded by the retained ledger. Selected original dependencies were inspected. The full external theories of analytic perturbation, elliptic regularity, stationary phase and nodal geometry were not all independently rebuilt.\par
\section*{Problem 1: the weak Euler calculation omits the hole pressure}
Location and claim. Theorem 1.3, p. 5; the definition of a stationary weak solution and the pressure formula, p. 33; and the first displayed calculation on p. 34, concluding its proof. The test class is explicitly every \(w\in C_c^\infty(\mathbb R^2,\mathbb R^2)\), not only divergence-free tests.\par
Write \(\Omega'\) for the bounded hole, \(\Gamma_1\) for the outer boundary and \(\Gamma_2=\partial\Omega'\) for the inner boundary. The source has \(u|_{\Gamma_j}=c_j\), \(\nabla u|_{\partial\Omega}=0\), and defines \[v=\begin{cases}(u_{x_2},-u_{x_1})&\text{in }\Omega,\\0&\text{outside }\Omega,\end{cases}\] together with the continuous pressure \[p=\begin{cases}
 0&\text{in the unbounded complementary component},\\
 -\frac12\bigl(|\nabla u|^2+\mu u^2-\mu c_1^2\bigr)&\text{in }\Omega,\\
 p_h:=-\frac{\mu}{2}(c_2^2-c_1^2)&\text{in }\Omega'.
 \end{cases}
 \tag{1}\] For \[I_D(w):=\int_D\bigl(v_i v_j\partial_iw_j+p\,\operatorname{div}w\bigr)\,\mathrm dx,\] the display on p. 34 asserts \[I_{\mathbb R^2}(w)=I_\Omega(w)
 =-\int_\Omega\bigl((v\cdot\nabla)v+\nabla p\bigr)\cdot w\,\mathrm dx=0.
 \tag{2}\] Both intermediate reductions in (2) fail for the defined pressure.\par
The pressure on the hole is nonzero. The radial Neumann branch used in the construction satisfies \[\psi''+\frac1r\psi'+\mu\psi=0,\qquad \psi'(a)=\psi'(1)=0.\] Its energy obeys the exact identity \[\frac{\mathrm d}{\mathrm dr}
 \left(\frac{(\psi')^2}{2}+\frac{\mu\psi^2}{2}\right)
 =-\frac{(\psi')^2}{r}.
 \tag{3}\] The eigenfunction is nonconstant and \(\mu>0\), so integration gives \(\psi(a)^2>\psi(1)^2\). Neither endpoint value vanishes, by uniqueness for the ODE. The second nonconstant radial eigenfunction has two interior zeros and therefore its endpoint values have the same sign.\par
The branch retains these radial boundary values: in the formula for \(\widetilde w_s\) on p. 31, both \(v_s\) and \(\overline\psi_{a(s)}'\) vanish on the boundary. Consequently \(c_2^2>c_1^2\) and \(p_h<0\). Alternatively, strict inequality at the radial point and continuity already suffice for small nonzero branch parameters. Thus the omitted term is not identically zero under the application’s hypotheses.\par
An admissible test and a decisive discrepancy. Choose a smooth compactly supported vector field \(w\) equal to \(x\) on a neighborhood of \(\overline{\Omega'}\) and zero near \(\Gamma_1\). Such a cutoff exists because the two boundary curves are disjoint compact sets. Then \[\int_{\Omega'}\operatorname{div}w\,\mathrm dx=2|\Omega'|,
 \qquad
 \int_{\Gamma_2}w\cdot\nu_\Omega\,\mathrm dS=-2|\Omega'|.
 \tag{4}\] The normals of \(\Omega\) and \(\Omega'\) on their common interface are opposite.\par
Inside \(\Omega\), the pointwise equation is correct. Direct differentiation gives \[(v\cdot\nabla)v+\nabla p
 =-(\Delta u+\mu u)\nabla u=0.
 \tag{5}\] Since \(v=0\) on both boundary components, the kinetic boundary flux vanishes. However, the pressure flux on \(\Gamma_2\) does not. Integration by parts therefore gives \[I_\Omega(w)=p_h\int_{\Gamma_2}w\cdot\nu_\Omega\,\mathrm dS
          =-2p_h|\Omega'|\ne0.
 \tag{6}\] The hole contributes \[I_{\Omega'}(w)=p_h\int_{\Omega'}\operatorname{div}w\,\mathrm dx
           =2p_h|\Omega'|.
 \tag{7}\] Thus the source’s assertion \(I_\Omega(w)=0\) is false, and the equality \(I_{\mathbb R^2}(w)=I_\Omega(w)\) is false as well. This is an exact flux calculation for an allowed test, not a numerical residual or a counterexample that violates the test class. The coordinate identity (5), the energy identity (3), and the cancellation in (6)–(7) were also checked with reviewer-written symbolic code.\par
Complete repair and counterevidence. The actual pair \((v,p)\) in (1) is a global weak solution. For an arbitrary test \(w\), the correct calculation is \[\begin{aligned}
 I_{\mathbb R^2}(w)
 &=I_\Omega(w)+p_h\int_{\Omega'}\operatorname{div}w\,\mathrm dx\\
 &=-\int_\Omega\bigl((v\cdot\nabla)v+\nabla p\bigr)\cdot w\,\mathrm dx\\
 &\quad+p_h\left(
   \int_{\Gamma_2}w\cdot\nu_\Omega\,\mathrm dS+
   \int_{\partial\Omega'}w\cdot\nu_{\Omega'}\,\mathrm dS\right)=0.
 \end{aligned}
 \tag{8}\] The outer pressure is zero and the two inner fluxes cancel. Equivalently, the stress tensor \(v\otimes v+pI\) has matching normal traces. In fact it is globally \(C^1\): from \(u\in C^2(\overline\Omega)\) and \(\nabla u=0\) on the boundary, one has \(\nabla p=-(D^2u+\mu uI)\nabla u\to0\), and the first derivatives of \(v_iv_j\) also tend to zero. Direct integration by parts on the whole plane therefore gives another complete repair. The pressure was already defined correctly in the source; it is the displayed integration argument that fails.\par
An interpretation using divergence-free tests would avoid the pressure integral, but it does not match the explicit quantifier on p. 33. An interpretation in which \(\Omega\) includes its hole contradicts the domain notation and the separate definition of \(p\) on \(\Omega'\). Neither reading makes (2) true as written.\par
Impact and classification. The invalid calculation is the step that turns the constructed annular eigenfunction into the whole-plane weak Euler solution promised in the title, abstract and Theorem 1.3. Under contract 0.7.1, this is a Major invalid derivation tied directly to an advertised core application. Its short repair and the survival of the actual construction do not change what the original calculation asserts. The finding does not invalidate the bifurcation theorem, the Pompeiu corollary, or the repaired Euler solution. No further hypothesis or new computation is needed for repair (8).\par
\section*{Problem 2: the rigidity proof changes its curvature weights}
Location and dependency. Theorem 6.1, p. 34; the parallel-curve relation (6.5), p. 37; the stationary-phase reduction (6.10), p. 39; and the definition of \(V\) leading to (6.11), p. 40. The theorem asserts that a smooth bounded planar domain admitting infinitely many orthogonal Neumann eigenfunctions locally constant on the boundary must be a disk or an annulus.\par
At the decisive point, the proof has distinct constants \(C_1,\ldots,C_m\), with \(C_1=0\), and a nonzero curvature \(k(s)\) on a small interval where \(1+C_\ell k(s)\ne0\). Equation (6.10) is \[\sum_{\ell=1}^m\widetilde\delta_\ell
 \left|\frac{k(s)}{1+C_\ell k(s)}\right|^{-1/2}=o(1).
 \tag{9}\] Taking a nonvanishing real or imaginary part and normalizing supplies unit vectors \(A_r\) and a subsequential limit \(A\) of norm one. The source then states that (9) can be rewritten using \[V(s)=\left(
 \frac{k(s)}{1+C_1k(s)},\ldots,
 \frac{k(s)}{1+C_mk(s)}
 \right)
 \tag{10}\] as \(A_r\cdot V(s)=o(1)\) and hence \(A\cdot V(s)=0\).\par
The actual vector determined by (9) is \[W(s)=\left(
 \left|\frac{k(s)}{1+C_1k(s)}\right|^{-1/2},\ldots,
 \left|\frac{k(s)}{1+C_mk(s)}\right|^{-1/2}
 \right).
 \tag{11}\] A linear relation among the components of \(W\) does not become the same linear relation among the components of \(V\).\par
Exact algebraic witness. Take regular values \[m=2,\qquad k=1,\qquad C_1=0,\quad C_2=3,\qquad
 A=\frac1{\sqrt5}(2,-1).\] All denominators are nonzero and the constants are distinct. Then \[W=(1,2),\qquad V=(1,\tfrac14),\qquad
 A\cdot W=0,\quad A\cdot V=\frac7{4\sqrt5}\ne0.
 \tag{12}\] Thus the proposed rewriting is not a linear-algebra identity, even at a single regular curvature value with a normalized nonzero coefficient vector. This is a counterexample to the displayed rewriting, not a constructed domain satisfying the spectral hypotheses of Theorem 6.1. The latter theorem is not claimed false.\par
A full repair of the curvature step. Retain \(W\) in the subsequence argument. The conclusion is then \(A\cdot W(s)=0\). Multiplying by \(\sqrt{|k(s)|}\) gives \[\sum_{\ell=1}^m A_\ell f_\ell(k(s))=0,\qquad
 f_\ell(k):=\sqrt{|1+C_\ell k|}.
 \tag{13}\] On a regular interval the signs of the linear factors are fixed. For any point \(k_0\) of that interval, \[f_\ell^{(j)}(k_0)
 =f_\ell(k_0)
   \left(\tfrac12\right)_{\!j}
   \left(\frac{C_\ell}{1+C_\ell k_0}\right)^j,
 \qquad 0\le j\le m-1,
 \tag{14}\] where \((\frac12)_j=\frac12(\frac12-1)\cdots(\frac12-j+1)\) is the falling factorial, with value one for \(j=0\).\par
Put \(q_\ell=C_\ell/(1+C_\ell k_0)\). The Wronskian determinant factors as \[\det\bigl[f_\ell^{(j)}(k_0)\bigr]_
       {\substack{0\le j<m\\1\le\ell\le m}}
 =
 \left(\prod_{\ell=1}^m f_\ell(k_0)\right)
 \left(\prod_{j=0}^{m-1}\left(\tfrac12\right)_{\!j}\right)
 \prod_{i<\ell}(q_\ell-q_i).
 \tag{15}\] Every factor is nonzero. Indeed, \[q_\ell-q_i=
 \frac{C_\ell-C_i}{(1+C_\ell k_0)(1+C_i k_0)}\ne0.\] This also includes \(C_1=0\); no division by \(C_1\) is used.\par
If \(k(s)\) were nonconstant, its continuous image would contain a nondegenerate interval. Equation (13) would give a fixed linear relation among the \(f_\ell\) on that interval, contradicting (15) because \(|A|=1\). Therefore \(k\) is constant. The subsequent circular-boundary and unique-continuation argument can then proceed with the intended conclusion.\par
Counterevidence and limits of the repair. The inverse-square-root factor in (9) is the stationary-phase amplitude and agrees with the preceding displayed formulas. Thus changing the definition of \(V\) to \(W\) is strongly supported by the source. Merely replacing the symbol in (10), however, does not prove the asserted noncontainment in a hyperplane; (13)–(15) provide the needed argument for the correct vector.\par
The common-phase calculation on p. 39 also contains local \(\xi/\eta\) and phase-sign inconsistencies. Correcting the derivative to \(p_1(s)\cdot\widetilde\eta(s)\) and absorbing the appropriate constant phase restores a common factor of modulus one. These are recorded with the proof clarification, not counted as independent Problems.\par
The stationary-phase grouping upstream of (9) was stress-tested separately. For each fixed \(s\), finitely many fixed imaginary shifts give a Vandermonde system in the distinct numbers \(\exp(b)\). Keeping the bounded coefficients \(\sigma_{n,\ell}\) until after this separation avoids an unsupported interchange with a growing exponential factor. This supplies counterevidence against an additional asserted failure upstream. The exact large-imaginary-shift theorem cited from Berenstein was not retrieved from the subscription-only primary source; the fixed-shift argument does not claim to certify that particular formulation.\par
Impact and classification. The claimed \(A\cdot V=0\) is the decisive input to curvature constancy and hence to Theorem 6.1. Its derivation from (9) is invalid. This is a Major defect in that rigidity derivation under contract 0.7.1, with the dependency explicit. The corrected Wronskian calculation restores this local step. It does not establish that every other imported part of the rigidity proof has been independently certified, and no failure of the main nonradial existence theorem follows from this issue.\par
\section*{Positive checks and bounded clarifications}
\subsection*{Crossing, nonresonance and the large-parameter calculation}
The derivative of a radial Neumann eigenfunction obeys the angular-mode-one Dirichlet equation. Sturm oscillation gives the index shift \[\mu_{0,n}(a)=\lambda_{1,n-1}(a),\qquad n\ge1.\] For fixed \(l\), the thin-annulus limits with \(h=(1-a)/\pi\) give \(\mu_{0,2}\sim4h^{-2}\) and \(\lambda_{l,0}\sim h^{-2}+l^2\). At the disk end, the relevant limits are \(j_{1,2}^2\) and \(j_{l,1}^2\), with the required strict inequality starting at \(l=4\). The official Bessel-zero interlacing formula [DLMF] and independent numerical values were inspected. Continuity supplies a crossing. The proof of nonresonance separately handles higher multiples of \(l\), higher radial indices, and the radial mode; it does not assume that the full unsymmetrized spectrum is simple.\par
The spectral checker reconstructed the first-order corrections and the second-order coefficients. With the source’s notation, the explicit functions \[\Phi_{10}(x)=(2x-\pi)\sin x+\frac32x(\pi-x)\cos x,
 \qquad
 \widetilde\Psi_{10}(x)=
 \left(\frac x2-\frac\pi4\right)\cos2x-\frac14\sin2x\] give \(\Lambda_{20}=-16\) and \(\widetilde\nu_{20}=3/4\). These are exact symbolic checks, distinct from numerical corroboration. The fixed-domain simple eigenbranches depend analytically on the common parameters, which supplies control of differentiated Taylor remainders. In particular the independent first-order calculation gives \[\mu'_{0,2}(a_l)-\lambda'_{l,0}(a_l)
 =\frac{2l^3}{\sqrt3\,\pi}+O(l^2)>0\] for sufficiently large \(l\). No invalid differentiation of a merely pointwise big-\(O\) estimate was established.\par
\subsection*{Functional setting, branch symmetry and analyticity}
The anisotropic space, following the framework in [FMW], controls an additional radial derivative. Its role is to keep the nonlinear pullback map and the boundary reconstruction in compatible spaces. The check traced the kernel, the weighted range orthogonality, the Fredholm index and the eigenvalue-derivative form of transversality through the bifurcation argument. The original Crandall–Rabinowitz statement [CR] was checked against the actual one-dimensional kernel and codimension-one range.\par
The first-order inner-radius expression in Theorem 1.1 also needs \(a'(0)=0\). This follows from a justified symmetry: rotation by \(\pi/l\) preserves the dihedral function space and changes the sign of its kernel vector. Using the kernel projection as the local parameter and an invariant complement, local uniqueness gives \(a(-s)=a(s)\). Thus the constant radial drift is \(o(s)\), and the stated leading cosine is consistent.\par
Spatial analyticity is distinct from analyticity in the bifurcation parameter. Near each boundary, \(u=c_i\ne0\) and \(\nabla u=0\); for \(z=u-c_i\) the equation is \(\Delta z+\mu z+\mu c_i=0\). The nonzero constant term supplies the nondegeneracy required by the inspected Kinderlehrer–Nirenberg free-boundary theorem [KN]. The source asserts a \(C^1\) parameter curve and analytic spatial boundaries, not an analytic parameter dependence.\par
\subsection*{Applications, rigidity and the numerical appendix}
The Pompeiu-type gluing has matching values and first derivatives across both interfaces. Its compactly supported function satisfies \[(\Delta+\mu)w=\mathbf{1}_{\Omega}-c\,\mathbf{1}_{\Omega'},
 \qquad c=c_2/c_1-1>0,\] in the distributional sense. Entire Helmholtz solutions therefore annihilate this source. The rigid-motion line has an inverse-convention slip, repaired by the standard change of variables; the universal integral identity survives.\par
For the Euler field, analyticity and nonconstancy imply that the gradient cannot vanish on an open subset, so \(\operatorname{supp}v=\overline\Omega\). The nonzero boundary values also explain the jump in the normal derivative of the velocity, while the velocity itself remains continuous. These properties are compatible with the weak, rather than \(C^1\), formulation. If the analytic field were radial on a nonempty open piece, its rotational derivative about that center would vanish analytically throughout the connected annular domain. The boundary would then lie on circles where the radial derivative vanishes. The noncircular branch therefore cannot be assembled from locally radial pieces.\par
The low-eigenvalue rigidity argument was checked separately from the infinite-spectrum theorem. It uses the derivatives generated by translations and rotations as Dirichlet eigenfunctions, eigenvalue comparison with multiplicities, and nodal topology. Its hypotheses and zero-based indexing were retained; no additional counterexample was established. The full cited Filonov/Friedlander and Reichel originals were not all retrieved; their role here is a checked statement and applicability interface, not a newly certified foundation.\par
Independent floating-point calculations for Appendix A give \[\begin{aligned}
 a_4&=0.1409893318207135179162\ldots,\\
 \mu_{0,2}(a_4)=\lambda_{4,0}(a_4)
     &=57.585155069669585459\ldots,\\
 \mu'_{0,2}(a_4)&=105.970673822067813524\ldots,\\
 \lambda'_{4,0}(a_4)&=0.1206696801319745464\ldots .
 \end{aligned}\] The reported derivative formulas agree with independent implicit differentiation and convergent centered differences. The numerical mode count gives 18 eigenvalues below the target, counting angular multiplicities, consistent with \(\mu_{18}\). Two printed index labels on p. 46 are transposed: the value \(44.0466317110\ldots\) belongs to \(\mu_{2,1}\) and \(64.1201718634\ldots\) to \(\mu_{3,1}\). Correcting those labels preserves the count.\par
The appendix itself distinguishes these computations from a computer-assisted proof. The present recomputation likewise supplies no certified interval enclosure. This limitation does not remove the separate analytic large-\(l\) result. Missing released author code is not treated as a mathematical defect.\par
\section*{Assessment and repair state}
@L0.15L0.11L0.27L0.17L0.21@ Item \& Assessment \& Repair verification \& Author adoption \& Remaining work
Item \& Assessment \& Repair verification \& Author adoption \& Remaining work
Problem 1 \& Major \& Exact hole/interface cancellation proves the intended global weak equation \& Pressure definition already correct; displayed proof still omits terms \& Replace the integration chain by (8)
Problem 2 \& Major \& Correct weights and nonzero Wronskian restore curvature constancy \& Correct amplitude appears in (6.10); vector below it remains inconsistent \& Correct the vector and supply its linear-independence argument
Large-\(l\) construction \& No failure established \& Spectral and bifurcation interfaces checked within scope \& Reviewed v2 only \& Imported analytic foundations and independent acceptance
Appendix \(l=4\) \& Numerical support \& Independent roots, derivatives and mode count agree \& Source labels numerical status \& Certified enclosures would be a separate validation task
Overall \& Initial review only \& Repairs and theorem validity kept distinct \& No verified subsequent correction \& Scientific decision and human review remain pending\par
\section*{Supporting checks}
The main source is https://arxiv.org/pdf/2309.07977v2, with its deposited TeX. The live arXiv record lists v2 as the final revised author manuscript. An additional 43-page author-site copy was retrieved and its identity page inspected; its bytes and 2023 PDF metadata differ from the assigned 50-page v2. It was used only for version comparison. The publisher landing-page check did not yield a usable publisher PDF. No claim is made that either finding remains present in the version of record.\par
All five assigned input records, totaling 1,034,853 bytes, were checked against actual bytes and hashes. The primary PDF SHA256 is\par
> 21e8e41c4f6944d2ca24550f334b1226b41cbed1c5d367e5f9e7c041ef1e4e19.\par
The TeX SHA256 is\par
> 99b1223737238a5756c35d52d373e91184dbdff56e4f3a77a4a580d8e29f55b5,\par
and the source-manifest SHA256 is\par
> 78cb1474c1a739cb63c95c56aadb7ac0a681aaf11e0edc01b2b37d1335de8ff8.\par
The evidence records retain exact source correspondence, computation inputs and outputs, calculation code, retrieved primary dependencies, and actual page observations. Search results did not identify a specific erratum or released paper code in the inspected records; this is neither an absence proof nor a claim of first discovery.\par
The actual native roles were `/root/resistance\_check` for the spectral proof, `/root/meanfield\_check` for the bifurcation and pullback analysis, and `/root/appendix\_check` for rigidity and Appendix A. They were bounded scientific assistants, not independent acceptance units. The main controller supplied the UI observation; this worker did not independently verify the UI or backend model.\par
The final report was genuinely compiled with the existing Tectonic 0.17.0 in untrusted mode. Every final report page and every used source-evidence page was actually viewed by the recorded reviewer. No author code, shell escape, installation or cloud export was used.\par
9 A. Enciso, A. J. Fernández, D. Ruiz and P. Sicbaldi, A Schiffer-type problem for annuli with applications to stationary planar Euler flows, arXiv:2309.07977v2, 12 August 2024. Reviewed author manuscript. M. G. Crandall and P. H. Rabinowitz, Bifurcation from simple eigenvalues, Journal of Functional Analysis 8 (1971), 321–340. Original scan consulted for the local bifurcation theorem. D. Kinderlehrer and L. Nirenberg, Regularity in free boundary problems, Annali della Scuola Normale Superiore di Pisa, Classe di Scienze (4) 4 (1977), 373–391. Original scan consulted for the analytic free-boundary interface. M. M. Fall, I. A. Minlend and T. Weth, The Schiffer problem on the cylinder and on the 2-sphere, arXiv:2303.17036v3, 13 May 2024. The relevant anisotropic-space and bifurcation interfaces were consulted; the complete paper is not independently certified here. NIST Digital Library of Mathematical Functions, Section 10.21, Equation 10.21.2a, https://dlmf.nist.gov/10.21. Official interlacing equation and its displayed image consulted.\par
\end{document}
